\section{Tesis 1}

\subsection{Título}
\textit{Exploring Emotion Recognition of Students in Virtual Reality Classrooms Through Convolutional Neural Networks and Transfer Learning Techniques} (Explorando el reconocimiento de emociones de los estudiantes en aulas de realidad virtual a través de redes neuronales convolucionales y técnicas de aprendizaje por transferencia) \cite{shomoye2024}.

\subsection{Autor}
Michael Shomoye, University of Calgary.

\subsection{Resumen}
``En los entornos educativos contemporáneos, comprender y evaluar el compromiso de los estudiantes a través de señales no verbales, especialmente las expresiones faciales, es fundamental. Tales señales han informado durante mucho tiempo a los educadores sobre los estados cognitivos y emocionales de los estudiantes, ayudándoles a adaptar sus métodos de enseñanza. Sin embargo, el auge de las plataformas de aprendizaje en línea y las tecnologías avanzadas como la Realidad Virtual (RV) desafían los modos convencionales de medir el compromiso de los estudiantes, especialmente cuando ciertas características faciales se oscurecen o están completamente ausentes. Este trabajo explora el potencial de las Redes Neuronales Convolucionales (CNN), específicamente un modelo de enfoque de entrenamiento personalizado adaptado de la arquitectura ResNet50, para reconocer y distinguir expresiones faciales sutiles en tiempo real, como la neutralidad, el aburrimiento, la felicidad y la confusión. La novedad del enfoque es doble: primero, se optimiza el poder de las CNN para analizar las expresiones faciales en plataformas de aprendizaje digital. Segundo, se innova para el contexto de la RV centrándose en la mitad inferior del rostro para abordar los desafíos de oclusión que plantea el uso de visores de RV. A través de una experimentación exhaustiva, se compara el rendimiento del modelo con la red neuronal residual 50 predeterminada (ResNet50) y se evalúa frente a conjuntos de datos de rostro completo y de rostro ocluido por RV. En última instancia, este esfuerzo tiene como objetivo proporcionar a los educadores una herramienta sofisticada para la evaluación en tiempo real del compromiso de los estudiantes en entornos de aprendizaje tecnológicamente avanzados, enriqueciendo la experiencia de enseñanza y aprendizaje.'' \cite{shomoye2024}

\subsection{Problema que busca resolver}
La tesis busca resolver la dificultad que enfrentan los educadores para evaluar el compromiso y el estado emocional de los estudiantes en entornos de aprendizaje virtuales y de Realidad Virtual. En un aula física, los profesores dependen de señales no verbales y expresiones faciales para ajustar su enseñanza, pero en entornos de RV, los visores ocultan la parte superior del rostro (oclusión), eliminando pistas visuales críticas y creando una desconexión entre la instrucción y las necesidades del estudiante \cite{shomoye2024}.

\subsection{Objetivos}

Objetivo General (inferido): Desarrollar un modelo de reconocimiento de emociones técnicamente robusto, pedagógicamente valioso y éticamente responsable que permita a los educadores obtener información sobre el compromiso de los estudiantes en tiempo real dentro de aulas de Realidad Virtual (RV) \cite{shomoye2024}.

Objetivos Específicos (inferidos):
\begin{enumerate}
    \item Comparar el rendimiento de la arquitectura ResNet50 con otros modelos de CNN prominentes como VGG19 y MobileNet en el contexto específico de reconocimiento facial en aulas de RV \cite{shomoye2024}.
    \item Evaluar la eficacia de un modelo personalizado basado en ResNet50 para diferenciar expresiones sutiles como neutralidad, aburrimiento, felicidad y confusión \cite{shomoye2024}.
    \item Determinar las técnicas de preprocesamiento óptimas (como el uso de imágenes en escala de grises) para mejorar el reconocimiento de emociones \cite{shomoye2024}.
    \item Superar el desafío de la oclusión facial causada por los visores de RV enfocándose en la parte inferior del rostro \cite{shomoye2024}.
\end{enumerate}

\subsection{Hipótesis}

La hipótesis no se encuentra explícitamente en el texto pero se infiere que un modelo de Red Neuronal Convolucional (basado en ResNet50) que utilice aprendizaje por transferencia y se centre en las características de la mitad inferior del rostro podría superar las limitaciones de oclusión de los visores de RV y clasificar con precisión estados emocionales clave del estudiante. Los resultados comprobaron esta hipótesis: el modelo personalizado no solo superó a la versión predeterminada de ResNet50 y a otros modelos, sino que también demostró una alta capacidad para distinguir el aburrimiento de la neutralidad tanto en rostros completos como ocluidos \cite{shomoye2024}.

\subsection{Resumen de la metodología empleada}
La investigación tomó una red neuronal preentrenada (ResNet50) y la adaptó para reconocer emociones faciales en el contexto específico de aulas de Realidad Virtual. Se usaron dos fuentes de imágenes: una base de datos pública de expresiones faciales (AffectNet) y un conjunto propio recolectado con 30 participantes en un entorno simulado. La metodología incluyó \cite{shomoye2024}:

\begin{itemize}
    \item \textbf{Preparación de imágenes:} Las fotografías se convirtieron a blanco y negro para simplificar el procesamiento. Luego se usó un programa de detección de puntos faciales para recortar el área del rostro, generando dos versiones de cada imagen: el rostro completo y únicamente la mitad inferior, para simular la obstrucción del visor de Realidad Virtual \cite{shomoye2024}.
    \item \textbf{Ajuste de la red neuronal:} Se desbloquearon las capas más profundas de la red preentrenada para que pudiera especializarse en reconocer emociones en este contexto. Se incorporaron técnicas para evitar que el modelo simplemente memorizara los ejemplos de entrenamiento en lugar de aprender a generalizar a imágenes nuevas \cite{shomoye2024}.
    \item \textbf{Entrenamiento:} Para compensar la cantidad limitada de imágenes disponibles, se generaron variaciones artificiales de cada imagen mediante volteos y rotaciones. Además, se usó un algoritmo que ajusta automáticamente la velocidad de aprendizaje de la red según su propio progreso \cite{shomoye2024}.
\end{itemize}

\subsection{Principales resultados}

Los principales hallazgos reportados por \cite{shomoye2024} son:

\begin{itemize}
    \item \textbf{Mejor arquitectura base:} De los tres tipos de redes neuronales comparados, ResNet50 mostró mayor precisión y mejor capacidad para generalizar a imágenes nuevas que las alternativas evaluadas (VGG19 y MobileNet), en el contexto específico de detectar emociones en entornos de Realidad Virtual \cite{shomoye2024}.
    \item \textbf{Mejora con el modelo ajustado:} La versión personalizada de la red neuronal, adaptada específicamente para esta tarea, superó de forma consistente a la versión estándar en todos los experimentos, incluyendo casos donde parte del rostro estaba cubierta por el visor de Realidad Virtual \cite{shomoye2024}.
    \item \textbf{Distinción entre aburrimiento y neutralidad:} El modelo logró diferenciar exitosamente estas dos emociones, que visualmente son muy similares, y mantuvo resultados aceptables incluso analizando únicamente la parte inferior del rostro \cite{shomoye2024}.
    \item \textbf{Ventaja de procesar imágenes en blanco y negro:} Convertir las imágenes a escala de grises antes de procesarlas redujo la carga computacional sin que el modelo perdiera capacidad para reconocer expresiones faciales \cite{shomoye2024}.
\end{itemize}

\subsection{Crítica de la tesis}
\begin{itemize}
    \item \textbf{Aspectos positivos:} La documentación alrededor de la metodología es clara y detallada, así como los diferentes experimentos realizados y sus resultados. También hay una clara relación entre las preguntas de investigación y las conclusiones. Finalmente, al ser un trabajo que involucra seres humanos, se destaca la consideración ética y el consentimiento informado. 
    \item \textbf{Aspectos por mejorar:} En primer lugar, los objetivos y la hipótesis no se enuncian explícitamente, lo que dificulta su identificación. En segundo lugar, la sección de conceptos introductorios es extensa y demasiado general. Por último, la discusión de las limitaciones del estudio es breve y no se observa mucha reflexión crítica sobre los posibles sesgos o limitaciones metodológicas.
\end{itemize}
 

